\section{Approach}
\label{sec_intro_approach}

The grouping module is initially built on a standard implementation of a GATv2~\cite{brody2022gatv2} model trained
with contrastive loss and $k$-means. This configuration is used as the baseline. The detected keypoints are used to 
build a graph per image
where each keypoint is a node and training forces same person keypoints closer together
in embedding space while separating different person ones. These embeddings are grouped using $k$-means 
to find person identities. The training pipeline uses a generated synthetic dataset 
with precise labels and can be further fine-tuned on COCO~\cite{lin2014coco}.
Three grouping heads are applied to the embeddings (to replace the presumed naive $k$-means): 
deep embedded clustering~\cite{xie2016dec}, slot attention~\cite{locatello2020slotattention}
and graph partitioning~\cite{yang2020gcnve} as well as 
ten variants of a novel skeleton-aware deep modularity network (SA-DMoN) which is an extension of
DMoN~\cite{tsitsulin2023dmon}. These all fail to outperform $k$-means, forcing a pivot
to focus on the embeddings themselves, leading to the development of the novel PG-GAT
which is a GATv2 with three skeleton-aware modifications to the attention mechanism. With the success  
of this module (as compared to the grouping attempts), its architecture, loss and fine-tune hyperparameters 
are then selected by Bayesian sweeps. Variants are then created to stress the design including a hyperbolic
embedding space, a triplet-based graph primitive, and visual-feature augmentation. Finally, 
a frozen-backbone person count head is added to remove the oracle-$K$ dependency, closing the 
loop with a fully autonomous end-to-end pipeline.

Evaluation covers three settings: the synthetic distribution, COCO with ground-truth
keypoints, and a full end-to-end pipeline where the keypoints come from a detector.
For the end-to-end setting that detector is a frozen HigherHRNet~\cite{cheng2020higherhrnet}.
Modularity and detector independence are then shown by applying the same grouping module,
without retraining, to detections from two further detector families:
YOLO11x-pose~\cite{ultralytics2024yolo11} and RTMO-L~\cite{lu2024rtmo}.